%% kapa: exam=1 round=26 year=2015 period=1 form=A subject=civil_law subject_ko=민법 seq=9 number=9
%% source: Q-Net 감정평가사 기출문제 (cst00309) artlSeq=5138001 file=extracted/1008/5138001/제26회_감정평가사 1차_기출/제26회_감정평가사 1차_기출/1교시_감평_A형.pdf
%% conversion: kordoc 4.24.51 -> tools/gen1.py (auto); review=proofread-sonnet; flags=none
\begin{question}{9}
부재자에 관한 설명으로 옳지 않은 것은? (다툼이 있으면 판례에 따름)
\begin{choices}
  \item 법원이 선임한 재산관리인이 부재자의 재산에 대해 보존행위를 함에는 법원의 허가를 얻어야 한다.
  \item 부재자가 재산관리인을 선임하였으나 부재자의 생사가 분명하지 않은 경우, 법원은 청구권자의 청구에 의하여 재산관리인을 개임할 수 있다.
  \item 재산관리인의 처분행위에 대한 법원의 허가는 장래의 처분행위뿐만 아니라 과거의 처분행위에 대한 추인을 위해서도 할 수 있다.
  \item 법원이 선임한 재산관리인이 부재자의 사망을 확인했더라도 법원에 의해 선임결정이 취소되지 않는 한 재산관리인은 계속하여 권한을 행사할 수 있다.
  \item 재산관리인은 보수청구권을 가지며, 재산관리로 인하여 과실 없이 입은 손해에 대해 배상을 청구할 수 있다.
\end{choices}
\end{question}
\answer{1}
